\section{Chapter 13 Refresher: A Human-Centric Model for Understanding LLM Errors --- Diagnostic Decomposition and Multi-Objective Reasoning}

\subsection*{Setting the Scene}
The closing chapter consolidates mathematical judgment rather than introducing a new heavy formalism. The field moved from single-score evaluation toward structured diagnosis because one scalar cannot describe heterogeneous failure causes. The critical insight is decomposition before intervention.

\subsection*{Notation Introduced in This Chapter}
\begin{itemize}[leftmargin=1.5em]
  \item Axis labels: structure, meaning, context, grounding, environment.
  \item Conditional slice: subset of cases defined by context/task conditions.
  \item Tradeoff surface: set of feasible points across accuracy, safety, latency, and cost.
\end{itemize}

\subsection*{What You Need To Recall Precisely}
\begin{itemize}[leftmargin=1.5em]
  \item \textbf{Failure-axis decomposition:} structure, meaning, context, grounding, environment are separate tests.
  \item \textbf{Conditional error analysis:} inspect slices and contexts, not only aggregate rates.
  \item \textbf{Correlation vs intervention:} observed association does not imply remediation leverage.
  \item \textbf{Multi-objective tradeoffs:} optimize over accuracy, safety, latency, and cost jointly.
  \item \textbf{Diagnostic ordering:} simple structural failures should be ruled out before attributing deep semantic or governance causes.
\end{itemize}

\subsection*{How These Results Build on Each Other}
\begin{enumerate}[leftmargin=1.5em]
  \item Observe failure instance and classify candidate axes.
  \item Check low-level structural and context conditions first.
  \item Test grounding and environment causes next.
  \item Select intervention matched to diagnosed failure pathway.
\end{enumerate}

\subsection*{Reminder Glossary (Term by Term)}
\begin{itemize}[leftmargin=1.5em]
  \item \textbf{Structure failure:} format/syntax/protocol violation.
  \item \textbf{Meaning failure:} semantic mismatch despite fluent output.
  \item \textbf{Context failure:} missing or misused prompt/task constraints.
  \item \textbf{Grounding failure:} unsupported factual claims.
  \item \textbf{Environment failure:} external system/tool mismatch.
\end{itemize}

\subsection*{Worked Mini Example (Axis Decision)}
Suppose an answer is fluent but cites a nonexistent document.
\begin{itemize}[leftmargin=1.5em]
  \item \textbf{Structure:} pass (syntax and format are valid).
  \item \textbf{Meaning:} partial pass (language coherent).
  \item \textbf{Grounding:} fail (claim cannot be verified against source evidence).
\end{itemize}
The corrective action should prioritize grounding controls (retrieval/citation verification), not merely prompt rephrasing.

\subsection*{Worked Example (Axis-Matched Intervention)}
Case: answer includes correct format and plausible language but wrong regulation citation.
\begin{itemize}[leftmargin=1.5em]
  \item Structure: pass.
  \item Meaning: partial pass.
  \item Grounding: fail (citation not verifiable).
\end{itemize}
Intervention path:
\begin{enumerate}[leftmargin=1.5em]
  \item add retrieval requirement from authoritative corpus,
  \item enforce citation-span verification before final response,
  \item reject answer if citation cannot be matched.
\end{enumerate}
Reasoning: this is not primarily a prompt-style failure; it is evidence-path failure.

\subsection*{Intuition Checkpoints}
\begin{itemize}[leftmargin=1.5em]
  \item Calibration and uncertainty: see Chapter 1 refresher.
  \item A plausible answer can still be ungrounded.
  \item One failure can belong to multiple axes; diagnosis is often multi-label.
  \item The strongest fix is the one that changes the relevant information pathway.
\end{itemize}

\subsection*{Further Refresh}
\begin{itemize}[leftmargin=1.5em]
  \item Evaluation slicing and retrieval grounding: see Chapter 10 refresher.
  \item Risk and assurance framing: see Chapter 11 refresher.
  \item Runtime constraints and failure budgets: see Chapter 12 refresher.
  \item \textbf{Murphy} --- probabilistic uncertainty reasoning.
  \item \textbf{Error analysis/evaluation notes} --- slicing and diagnostic methodology.
\end{itemize}
